% Background chapter body. Included by thesis.tex under \chapter{BACKGROUND}.

\section{AI Coding Assistants and the Agentic Model}
\label{sec:coding-assistants}

How developers work with AI has changed quickly, from tools that completed a line of code to systems that carry out whole programming tasks on their own. The earliest of these were predictive. Integrated development environments have offered identifier and interface completion since the 1990s, and by the late 2010s those suggestions were ranked by machine-learning models trained on large code corpora. Such systems retrieve and rank symbols that already exist; they do not compose new code. The generative shift came when large language models (LLMs) trained on code began producing novel, multi-line completions inline in the editor, the form that GitHub Copilot brought to a wide audience in 2021.\footnote{GitHub Copilot technical preview, June 2021: \url{https://github.blog/2021-06-29-introducing-github-copilot-ai-pair-programmer/}.} A second and distinct mode of interaction followed in late 2022, when general-purpose conversational assistants reached the public. While these were not explicitly coding tools, developers used them broadly, to explain unfamiliar code, stand in for documentation, or think through a problem as much as to draft a snippet to paste back. This conversational mode was soon folded into the editor, and then extended past advice into action. The assistant started editing files, running commands, and revising its work against the results. Tools of this last kind, which act on a project rather than only commenting on it, are called \emph{agentic}; they emerged over 2024 and 2025.\footnote{Early examples include Devin, announced in March 2024 as an autonomous ``AI software engineer'' acting on a project from its own sandboxed environment rather than the developer's editor (\url{https://www.cognition.ai/blog/introducing-devin}), and GitHub Copilot's agent mode, introduced in February 2025, which works inside the editor to plan a task, edit across multiple files, run commands, and iterate on the results (\url{https://code.visualstudio.com/blogs/2025/02/24/introducing-copilot-agent-mode}).} It is this conversational line, rather than inline completion, that the agentic systems studied in this thesis descend from.

We use the term \emph{agent} in a specific sense throughout this thesis. In its classical formulation, an agent is anything that perceives its environment through sensors and acts upon that environment through actuators \cite{russell2021aima}. Brought into the setting of LLMs, an agent is built around an LLM, with the ability to perceive and use external tools and resources, extending it beyond producing text into planning, memory, and action \cite{liu2024llmagents}. What distinguishes such an agent from a fixed program is where control of the process sits. An agent directs its own steps and tool use, choosing what to do next, whereas a workflow drives the model along predefined paths \cite{anthropic2024agents}. The frameworks evaluated in this thesis are agentic in precisely this sense, the model rather than a fixed script driving the work.

The effect of these tools on developer productivity is measured, and contested, and the evidence spans different generations of the technology. Field experiments with an earlier generation of assistant reported substantial gains. In three randomised controlled trials run at Microsoft, Accenture, and a Fortune-100 company, several thousand professional developers given an AI coding assistant completed around a quarter more tasks, with the clearest gains among the less experienced \cite{cui2026genai}. Evidence from agentic tools points the opposite way. In a trial run by METR in early 2025, with the Cursor editor\footnote{Cursor, an AI code editor: \url{https://cursor.com}.} driving Claude 3.5 and 3.7 Sonnet, experienced open-source maintainers working on their own repositories were left roughly 19\% slower, even as those same developers expected to be sped up \cite{becker2025metr}. The two studies differ in their tools, their tasks, and the developers they observed, so the effect reads as conditional rather than settled. Both, like the wider productivity literature, ask only whether AI assistance helps; neither varies the workflow or framework through which a developer works with the agent, which is the dimension this thesis sets out to compare.

\section{Context Windows and Context Decay}
\label{sec:context-windows}

Whatever its form, an agent operates within a context window: the bounded span of tokens the underlying model can attend to at once. Into it go the system prompt, the instructions the agent was given, the files and tool output it has read, and its own earlier turns; a long task steadily fills it. A frontier model's context window now spans several hundred thousand to a million tokens, yet even that can be too small to hold an entire repository and run a complex task end to end.

A model does not even use the context it has evenly. Controlled experiments on long-context models found a U-shaped pattern in which performance is highest when the relevant information sits at the start or the end of the input and falls sharply when it lands in the middle, even for models built for long contexts \cite{nliu2023lostmiddle}. Usable capacity is therefore smaller than the raw token count suggests.

For an agentic coding session the consequence is direct. As the session runs, every file read, tool called, and command executed consumes context. As that context fills, the agent is more likely to forget earlier decisions, lose track of completed work, or drift from the original goal \cite{liu2024llmagents}. Summarising earlier context or retrieving it on demand buys room, but only by discarding detail, trading fidelity for space.

This accumulation also carries a direct monetary cost. Each step is produced by running the model over the entire session so far, supplied to it as input, and a call is priced by the number of tokens it processes. Every step therefore costs more than the one before it, and the cumulative cost of a long task grows faster than the number of steps it takes \cite{gao2025turncontrol}.

\placeholder{Forthcoming. Spec-driven development and its single-agent
reference implementation; single- versus multi-agent coordination; and a
compact positioning of the gap this thesis addresses.}
